\section{The un-retracted iterate as a running weighted average}
\label{sec:running_average}

\begin{lemma}[Running-average representation of the un-retracted iterate]
\label{lem:running_average}
Consider the SGDM dynamics of \Cref{def:dynamics} with step
$\alpha\in(0,1]$ and the convention $V_0\coloneqq x_0$. For every horizon
$T\ge1$ the un-retracted iterate admits the convex-combination
representation
\begin{align}\label{eq:running_average}
   \tilde x_T \;=\; \sum_{k=0}^{T} w_{T,k}\, V_k,
   \qquad
   w_{T,k}\ge0,
   \qquad
   \sum_{k=0}^{T} w_{T,k} = 1,
\end{align}
where the weights are the attention (softmax) weights
$w_{T,k}=e^{s_{T,k}}/\sum_{j=0}^{T}e^{s_{T,j}}$ of the head, with scores
$s_{T,k}\coloneqq\inner{q}{k_k}$ satisfying $\abs{s_{T,k}}\le S$
(\Cref{ass:bounded_architecture}(4)). Consequently the on-sphere iterate is
the exact rescaling of this convex combination,
\begin{align}\label{eq:onsphere_direction}
   x_T \;=\; \sqrt d\,\frac{\tilde x_T}{\norm{\tilde x_T}_2}
   \;=\; \sqrt d\,\frac{\sum_{k=0}^{T}w_{T,k}V_k}{\bigl\|\sum_{k=0}^{T}w_{T,k}V_k\bigr\|_2}
   \qquad (\text{assuming }\tilde x_T\neq0),
\end{align}
the initial weight obeys $w_{T,0}\le\max_{0\le k\le T}w_{T,k}\le e^{2S}/T$
(so $w_{T,0}=O(1/T)$), and $\tau\coloneqq T^{-1}\sum_{t=1}^{T}w_{T,t}\le1/T$.
The convex-combination identity $\sum_k w_{T,k}=1$ holds for the
un-retracted iterate $\tilde x_T$ only.
\end{lemma}

\begin{proof}
The proof has two parts: the softmax representation (an algebraic identity)
and the initial-weight decay (a consequence of \Cref{lem:quad_variation}).

\smallskip\noindent\textbf{Softmax-average representation.} A single
attention head at the decoding position computes a softmax-weighted average
of the value vectors it attends to, including the residual self-term
$V_0=x_0$ (\cite{vaswani2017attention}): with scores
$s_{T,k}=\inner{q}{k_k}$ for $k=0,\dots,T$,
\begin{align*}
   \tilde x_T \;=\; \sum_{k=0}^{T} w_{T,k}\,V_k,
   \qquad
   w_{T,k} \;=\; \frac{e^{s_{T,k}}}{\sum_{j=0}^{T} e^{s_{T,j}}}.
\end{align*}
Each $w_{T,k}>0$ as a ratio of positive exponentials, and
$\sum_{k=0}^{T} w_{T,k}=1$ because the numerators sum to the denominator;
this is \Cref{eq:running_average}. The EMA recurrence
$\tilde x_t=(1-\alpha)x_{t-1}+\alpha V_t$ of \Cref{def:dynamics} is the
recency specialisation in which the score increments are constant
($s_{T,k}-s_{T,k-1}=\log\frac{1-\alpha}{\alpha}\cdot(\cdots)$); the
representation above is exact and does not depend on that specialisation.
Since $\tilde x_T$ is a nonnegative combination of the $V_k$ and
$\sum_k w_{T,k}=1$, the on-sphere retraction
$x_T=\sqrt d\,\tilde x_T/\norm{\tilde x_T}_2$ is the exact rescaling
\Cref{eq:onsphere_direction} (the scalar $\norm{\tilde x_T}_2$ is the only
thing the retraction changes; the direction is that of the convex
combination).

\smallskip\noindent\textbf{Initial-weight decay.} The weight $w_{T,0}$ is
one of the $T+1$ softmax weights, hence is at most the maximal weight,
which \Cref{lem:quad_variation} (applied with the same bounded scores
$\abs{s_{T,k}}\le S$, now over the $T+1$ indices $k=0,\dots,T$) bounds by
$\max_{0\le k\le T}w_{T,k}\le e^{2S}/(T+1)\le e^{2S}/T$. Hence
$w_{T,0}=O(1/T)$, and the initial state washes out as $T\to\infty$. The
average-weight bound $\tau=T^{-1}\sum_{t=1}^{T}w_{T,t}\le T^{-1}\cdot1=1/T$
is immediate from $\sum_t w_{T,t}\le1$. This completes the proof.
\end{proof}

\begin{remark}[The identity is for $\tilde x_T$, never $x_T$]
\label{rem:running_average_untransformed}
\Cref{lem:running_average} is the \emph{only} source of the
convex-combination identity in the paper, and it applies strictly to the
un-retracted iterate $\tilde x_T$. The retracted iterate
$x_T=\sqrt d\,\tilde x_T/\norm{\tilde x_T}$ does \emph{not} inherit
$\sum_k w_{T,k}=1$ (it has norm $\sqrt d$, not the norm of a convex
combination of the $V_k$); the passage from $\tilde x_T$ to $x_T$ is the
exact rescaling \Cref{eq:onsphere_direction}, whose effect on a fixed-row
logit is the work of \Cref{lem:retraction_stability}. Downstream,
\Cref{lem:signal_floor} and \Cref{lem:incorrect_max} establish their bounds
on $\tilde x_T$ using \Cref{eq:running_average}, and only then transfer to
$x_T$.
\end{remark}
